%!TEX root = ../main.tex
\chapter*{Sommario}
\addcontentsline{toc}{chapter}{Sommario}

Il controllo semaforico su una rete di intersezioni è un problema di
coordinamento distribuito: la fase scelta a un incrocio influenza, con
pochi secondi di ritardo, il traffico che raggiunge gli incroci vicini. Le
soluzioni classiche, a tempi fissi o reattive sulla sola pressione locale
(MaxPressure), non tengono conto di questa dipendenza spaziale in modo
esplicito. MetaSTGAT (Wang et al., 2022) affronta il problema con un agente
di reinforcement learning per intersezione, coordinato da una \emph{graph
attention network} i cui pesi sono generati dinamicamente da un modulo di
meta-learning condizionato su caratteristiche locali della rete stradale.

Questo lavoro replica da zero l'architettura e la procedura di training di
MetaSTGAT su un ambiente costruito su \textsc{CityFlow}, con reti a griglia
di tre dimensioni ($4\times4$, $5\times5$, $6\times6$) e traffico generato
parametricamente: ogni configurazione include un'arteria a densità
artificialmente maggiorata, usata per misurare l'equità del controllo tra
direzioni di traffico diverse.
Ogni differenza introdotta rispetto alla procedura originale, dalla
ricompensa alla struttura del replay buffer, è isolata da un sistema di
preset di ablation e motivata singolarmente invece di essere presentata come
un insieme indistinto di miglioramenti.

I risultati sperimentali confrontano le baseline classiche (Fixed-Time,
MaxPressure) con tre varianti del modello Pro che differiscono per il peso
$\alpha$ della penalità anti-starvation ($\alpha \in \{0.08, 0.04, 0.00\}$),
con lo studio di ablation completo sui sei preset previsti e con un
confronto tra tre meccanismi di comunicazione spaziale a parità di resto
dell'architettura: il \emph{graph attention} meta-condizionato del modello di
riferimento (Meta-GAT), una convoluzione grafica meta-condizionata
(Meta-GCN) e una propagazione ondulatoria meta-condizionata ispirata a SONAR
(Meta-SONAR). Rimuovere del tutto la penalità anti-starvation ($\alpha=0$)
fa collassare il max wait verso i valori di MaxPressure e allo
stesso tempo peggiora il tempo di viaggio massimo rispetto alle varianti con
$\alpha>0$, escludendo un semplice compromesso tra equità e throughput. Lo
studio di ablation isola un effetto analogo su scala più ampia: disattivare
singolarmente i meccanismi dell'algoritmo di apprendimento (Double DQN, il
replay prioritizzato, la componente ricorrente della Meta-LSTM) lascia il
modello vicino al riferimento su ogni metrica, mentre rimuovere insieme il
vincolo di visibilità, il segnale di attesa nello stato e la ricompensa
riscritta riporta il comportamento vicino a quello di un controllore
reattivo privo di garanzie di equità.

Il confronto tra meccanismi spaziali mostra un divario netto tra la
convoluzione grafica, priva di un meccanismo di compatibilità appreso tra
nodi, e gli altri due: Meta-GCN resta indietro su ogni metrica, mentre
Meta-GAT e Meta-SONAR restano sostanzialmente alla pari, con il modello di
riferimento a un solo layer che si conferma il più equo dell'intero
confronto. La rappresentazione del tempo di attesa nello stato è infine
un'ultima scelta di design verificata sperimentalmente: sostituire la
normalizzazione lineare, che rende indistinguibili tutte le attese oltre i
100 secondi, con una funzione $\tanh$ che continua a crescere ben oltre
quella soglia migliora ogni metrica riportata, motivo per cui questa
rappresentazione è adottata per l'intero parco di modelli di questo lavoro,
non solo per la configurazione su cui è stata isolata.
